% =====================================================================
%  FICHA CATALOGRÁFICA — MODELO  (NBR 14724:2024, 4.2.1.1.2)
% ---------------------------------------------------------------------
%  ATENÇÃO: a ficha DEFINITIVA deve ser elaborada por bibliotecário(a) da
%  UNIVALI (Código de Catalogação vigente). Este é apenas um modelo
%  preenchível; os campos entre <...> são fornecidos pela Biblioteca.
%  Elemento OBRIGATÓRIO apenas na versão final depositada. Para incluir,
%  descomente em thesis.tex:
%      \incluifichacatalografica{tex/pre-textual/ficha-catalografica.tex}
% =====================================================================

\begin{center}
\setlength{\fboxsep}{1.4em}%
\fbox{%
  \begin{minipage}{12.5cm}
    \fontsize{10pt}{13pt}\selectfont
    \setlength{\parindent}{0pt}%
    \begin{tabular}{@{}p{1.7cm}@{\hspace{0.3cm}}p{10.0cm}@{}}
      & <Sobrenome, Nome>. \\[2pt]
      <F000a> &
        <Título do trabalho: subtítulo>\,/\,<Nome Completo do Autor>. --
        \imprimirlocal, <ano>. \\[2pt]
      & \mbox{}\hspace{0.5cm}<000>~f. : il. ; 30~cm. \\[8pt]
      & \imprimirtipotrabalho{} (\tipodocurso{} <em Área>) -- <Instituição>,
        \imprimirlocal, <ano>. \\[2pt]
      & \mbox{}\hspace{0.5cm}Orientador: \titulacaoorientador{} \imprimirorientador. \\[8pt]
      & 1.~<Palavra-chave>. 2.~<Palavra-chave>. 3.~<Palavra-chave>.
        I.~<Sobrenome do orientador, Nome>. II.~Título. \\[8pt]
      & \hfill CDD: <000> \\
    \end{tabular}
  \end{minipage}%
}

\vspace{6pt}
{\fontsize{9pt}{11pt}\selectfont
 Ficha catalográfica elaborada por <Nome do(a) bibliotecário(a)> -- CRB <00/0000>\par}
\end{center}
